\chapter{Implementacja}
\label{chap:implementacja}

\section{Struktura kodu}
\label{sec:struktura-kodu}

Implementacja aplikacji została podzielona zgodnie z podziałem odpowiedzialności opisanym w \cref{chap:projekt-programu}. Kod źródłowy obejmuje moduły związane z interfejsem użytkownika, parserem, stemplami elementów, macierzami oraz solverem. Taki układ odpowiada naturalnemu przepływowi danych: od tekstowego opisu obwodu, przez reprezentację elementów, aż do układu równań rozwiązywanego numerycznie i prezentacji wyniku.

Warstwa widoku odpowiada za przyjmowanie netlisty od użytkownika i wybór formy prezentacji wyniku. Warstwa parsera zawiera typy odpowiedzialne za rozpoznawanie elementów i komend analizy. Warstwa stempli zawiera klasy opisujące konkretne elementy obwodu. Warstwa matematyczna udostępnia macierze oraz operacje potrzebne do ich składania i rozwiązywania. Solver łączy część obliczeniową, ponieważ otrzymuje gotową listę stempli i wykonuje analizę wskazaną przez komendę.

W implementacji przyjęto założenie, że elementy nie powinny samodzielnie interpretować tekstu wejściowego. Każdy konstruktor elementu otrzymuje już przetworzone wartości liczbowe: numery węzłów, parametry oraz ewentualne indeksy dodatkowych równań. Dzięki temu błędy składniowe pozostają w parserze, a błędy fizycznie lub numerycznie niepoprawnych parametrów pozostają w warstwie stempli.

\section{Parser i wynik parsowania}
\label{sec:implementacja-parser}

Parser pracuje na tekście zakodowanym jako UTF-8. Tekst może pochodzić z pliku albo bezpośrednio z pola edycji w interfejsie użytkownika. Każda linia jest dzielona na tokeny, a pierwszy token decyduje o dalszej interpretacji. Jeżeli token odpowiada komendzie analizy, parser zapisuje ją jako polecenie dla solvera. Jeżeli token odpowiada elementowi, parser sprawdza liczbę parametrów oraz zapamiętuje linię do późniejszego utworzenia stempla.

Rezultatem parsowania jest struktura zawierająca listę stempli, komendę analizy oraz listę węzłów przeznaczonych do wizualizacji. Taki wynik jest wygodny, ponieważ stanowi pełne wejście dla solvera i widoku wyniku, ale nie wiąże parsera z konkretną metodą obliczeń.

\begin{lstlisting}[language=Swift,caption={Model wyniku parsowania},label={lst:parser-result}]
struct ParserResult {
    let stamps: [Stamp]
    let command: SolverCommand
    let showNodes: [Int]
}
\end{lstlisting}

Ważnym elementem parsera jest typ \texttt{ElementType}. Odpowiada on za powiązanie oznaczeń tekstowych z obsługiwanymi elementami programu. Ten typ przechowuje również informacje o wymaganej liczbie tokenów oraz o pozycjach tokenów, które należy traktować jako numery węzłów.

\begin{lstlisting}[language=Swift,caption={Uproszczona rola typu ElementType},label={lst:element-type}]
enum ElementType {
    case resistor
    case capacitor
    case DCCS
    case DCVS
    case ACVS
    case CCCS
    case VCCS
    case CCVS
    case VCVS
    case BJT
}
\end{lstlisting}

Takie rozwiązanie ogranicza liczbę miejsc, w których trzeba zmieniać kod po dodaniu kolejnego elementu. Nowy element wymaga rozszerzenia słownika typów, określenia liczby parametrów oraz dodania funkcji tworzącej właściwy stempel.

\section{Tworzenie stempli}
\label{sec:implementacja-stemple}

Podstawowym typem wspólnym dla elementów jest \texttt{Stamp}. Udostępnia on metody zwracające lokalną macierz przewodności oraz lokalny wektor wymuszeń. Klasy konkretnych elementów nadpisują te metody zgodnie z równaniami wynikającymi z metody ZMPW.

\begin{lstlisting}[language=Swift,caption={Wspólny interfejs stempla},label={lst:stamp-interface}]
class Stamp {
    func getGMatrix(_ context: StampContext) throws -> Matrix? {
        return nil
    }

    func getIMatrix(_ context: StampContext) throws -> Matrix? {
        return nil
    }
}
\end{lstlisting}

Argument \texttt{StampContext} przekazuje do elementu informacje potrzebne w danym rodzaju analizy. Dla analizy czasowej zawiera między innymi aktualny czas, krok czasowy, bieżące przybliżenie rozwiązania oraz rozwiązanie z poprzedniego kroku. Dzięki temu ten sam interfejs stempla może obsłużyć elementy liniowe, elementy dynamiczne i elementy nieliniowe.

Dla elementów pasywnych, takich jak rezystor, najważniejsza jest macierz \texttt{G}. Dla niezależnego źródła prądowego zasadniczy wkład trafia do wektora \texttt{I}. Źródła napięciowe i wybrane źródła sterowane wymagają dodatkowego równania, dlatego ich stemple są większe i obejmują dodatkowe wiersze oraz kolumny.

Z punktu widzenia implementacji istotne jest to, że stempel nie musi znać pełnego rozmiaru układu. Zwraca macierz o takim rozmiarze, który jest potrzebny do zapisania jego największego indeksu. Dopiero solver tworzy macierz globalną o rozmiarze dopasowanym do największego stempla w całej liście.

\section{Macierze}
\label{sec:implementacja-macierze}

Macierz jest reprezentowana przez liczbę wierszy, liczbę kolumn oraz jednowymiarową tablicę wartości typu \texttt{Double}. Dostęp do elementu macierzy odbywa się przez przeliczenie indeksu wiersza i kolumny na indeks tablicy. W implementacji numery węzłów i równań są zgodne z notacją obwodową, dlatego indeks \texttt{0} oznacza masę i nie jest zapisywany w macierzy.

Dodawanie wartości do macierzy ma charakter akumulacyjny. Jeżeli dwa elementy wnoszą wkład do tej samej pozycji, wartości są sumowane. Jest to bezpośrednie odzwierciedlenie składania równań ZMPW, gdzie każdy element dopisuje lokalny fragment do globalnego układu.

\begin{lstlisting}[language=Swift,caption={Zasada dodawania wartości do macierzy},label={lst:matrix-add}]
func add(_ row: Int, _ column: Int, _ value: Double) throws {
    if row == 0 || column == 0 || value == 0 { return }
    values[columns * (row - 1) + (column - 1)] += value
}
\end{lstlisting}

W programie wyróżniono klasy pomocnicze \texttt{GMatrix}, \texttt{IMatrix} oraz \texttt{VMatrix}. Pierwsza reprezentuje macierz układu, druga wektor prawej strony, a trzecia wektor rozwiązania. Ten podział nie zmienia sposobu przechowywania danych, ale zwiększa czytelność kodu.

\section{Solver}
\label{sec:implementacja-solver}

Solver rozpoczyna pracę od wyznaczenia wymaganego rozmiaru macierzy globalnej. W tym celu przegląda wszystkie stemple i sprawdza rozmiary zwracanych przez nie macierzy. Następnie tworzy pustą macierz \texttt{G} oraz pusty wektor \texttt{I}, po czym dodaje do nich wkład każdego elementu.

\begin{lstlisting}[language=Swift,caption={Schemat składania macierzy globalnej},label={lst:solver-assembly}]
for stamp in stamps {
    try gMatrix.add(stamp.getGMatrix(context))
    try iMatrix.add(stamp.getIMatrix(context))
}
\end{lstlisting}

Dla analizy punktu pracy układ ma postać
\begin{equation}
  \mathbf{G}\mathbf{x} = \mathbf{I}.
\end{equation}
Po złożeniu macierzy solver rozwiązuje ten układ i zwraca wektor niewiadomych. Wektor zawiera napięcia węzłowe oraz dodatkowe niewiadome wprowadzone przez źródła napięciowe i źródła sterowane prądem.

Analiza czasowa korzysta z tej samej struktury organizacyjnej. Różnica polega na tym, że macierze dla elementów dynamicznych zależą od kroku czasowego i poprzednich stanów układu. Przed rozpoczęciem analizy \texttt{.tran} solver wyznacza punkt pracy, a następnie traktuje go jako warunek początkowy dla pierwszego kroku czasowego.

W przypadku elementów nieliniowych, takich jak tranzystor bipolarny, solver używa iteracji Newtona-Raphsona. W każdej iteracji stemple są obliczane dla aktualnego przybliżenia rozwiązania, a następnie rozwiązywany jest zlinearyzowany układ równań. Proces kończy się po osiągnięciu zadanej tolerancji albo po przekroczeniu maksymalnej liczby iteracji. W aktualnej implementacji limit ustawiono na 150 iteracji, co pozwala stabilnie policzyć testowy układ wtórnika emiterowego.

\section{Interfejs użytkownika i wizualizacja wyników}
\label{sec:implementacja-wizualizacja}

Interfejs użytkownika został zaimplementowany w SwiftUI. Główny widok ma trzy podstawowe stany: edycję opisu układu, prezentację wyniku punktu pracy oraz prezentację wyniku analizy czasowej. W stanie edycji użytkownik może wpisać netlistę ręcznie albo przeciągnąć plik \texttt{.txt}; po upuszczeniu pliku jego zawartość zastępuje tekst w edytorze.

Po naciśnięciu przycisku uruchomienia program przekazuje tekst do parsera. Jeżeli komendą analizy jest \texttt{.op}, widok pokazuje wartości kolejnych niewiadomych wektora rozwiązania. Jeżeli komendą jest \texttt{.tran}, solver zwraca kolejne próbki czasowe, a widok przechodzi do wykresu. W obu przypadkach użytkownik może wrócić do edytora bez utraty wpisanego opisu układu.

Do prezentacji wyników analizy czasowej dodano moduł \texttt{Plotter}. Nie wykonuje on obliczeń obwodowych, lecz przekształca wynik solvera na serie punktów opisane czasem i napięciem wybranego węzła. Lista węzłów pochodzi z komendy \texttt{.show}, dzięki czemu opis wejściowy decyduje o tym, które przebiegi zostaną pokazane.

Warstwa widoku korzysta z przygotowanych serii i rysuje wykres napięcia w funkcji czasu. Takie rozdzielenie ułatwia dalszy rozwój programu: solver pozostaje odpowiedzialny wyłącznie za obliczenia, a wizualizacja może zostać rozbudowana niezależnie, na przykład o eksport danych albo porównywanie kilku symulacji.

\section{Obsługa błędów w implementacji}
\label{sec:implementacja-bledy}

Błędy zostały podzielone na kilka grup. \texttt{ParserError} opisuje problemy z plikiem wejściowym, takie jak nieznany element, zła liczba parametrów, niepoprawny typ parametru, brak komendy, więcej niż jedna komenda analizy, więcej niż jedna komenda \texttt{.show} albo użycie \texttt{.show} bez analizy czasowej. Każdy błąd parsera związany z konkretną linią zawiera jej numer.

\texttt{StampParameterError} opisuje problemy wykryte przez elementy. Przykładem jest ujemny indeks węzła albo wartość parametru, której nie można użyć w obliczeniach. Parser przechwytuje takie błędy i zamienia je na informację diagnostyczną odnoszącą się do linii pliku. Dzięki temu użytkownik otrzymuje komunikat związany z wejściem, a kod elementów nie musi znać kontekstu tekstowego.

Osobną grupę tworzą błędy solvera i macierzy. Solver zgłasza między innymi pustą listę stempli, nieobsługiwaną komendę, niepoprawne parametry analizy czasowej, dywergencję numeryczną oraz brak zbieżności Newtona-Raphsona. Błędy macierzy obejmują na przykład przekroczenie rozmiaru macierzy, niezgodne wymiary albo błąd rozwiązania układu.

\section{Uwagi implementacyjne}
\label{sec:uwagi-implementacyjne}

Implementacja w języku Swift pozwala korzystać z typów statycznych, wyjątków oraz bibliotek systemowych dostępnych na platformie macOS. Dla małych macierzy używana jest własna implementacja rozwiązania układu, natomiast dla większych macierzy można wykorzystać funkcje biblioteki LAPACK dostępne przez framework Accelerate. Takie rozwiązanie pozwala zachować prostotę w przypadkach testowych i jednocześnie pozostawia drogę do obliczeń na większych układach.

Najważniejszą cechą implementacji jest zachowanie prostego kontraktu między warstwami. Parser zwraca stemple i komendę, stemple zwracają macierze lokalne, macierze potrafią dodawać wkłady, a solver wykonuje obliczenia. Dzięki temu każda część programu może być rozwijana i testowana niezależnie.
